% Part: lambda-calculus
% Chapter: lambda-definability
% Section: truth-values

\documentclass[../../../include/open-logic-section]{subfiles}

\begin{document}

\olfileid{lam}{ldf}{tvr}
\olsection{Nilai Bebeneran lan Relasi}

Nilai bebeneran bisa dikodeake
ing kalkulus lambda murni
kaya mangkene:
\begin{align*}
  \fn{true} & \ident \lambd[x][\lambd[y][x]]\\
  \fn{false} & \ident \lambd[x][\lambd[y][y]]
\end{align*}

Nilai bebeneran diwujudake
minangka \emph{selektor},
yaiku fungsi kang nampa
rong argumen lan ngasilake
salah sijine.
Nilai $\fn{true}$ milih
argumen kapisan, dene
$\fn{false}$ milih
argumen kapindho.
Contone, $\fn{true}\, M N$
mesthi direduksi dadi $M$,
dene $\fn{false}\, M N$
mesthi direduksi dadi~$N$.

\begin{defn}
Relasi $R \subseteq \Nat^k$
diarani !!{lambda definable}
yen ana suku~$R$ kang nyukupi
\begin{align*}
  R\, \num{n_1} \dots \num{n_k} & \bred \fn{true}
  \intertext{nalika $R(n_1, \dots, n_k)$ bener, lan}
  R\, \num{n_1} \dots \num{n_k} & \bred \fn{false}
\end{align*}
yen ora.
\end{defn}

Contone, relasi
$\fn{IsZero} = \{0\}$
kang bener tumrap $0$
lan mung $0$ iku
!!{lambda definable}
dening
\[
  \fn{IsZero} \ident \lambd[n][n (\lambd[x][\fn{false}])\, \fn{true}].
\]
Kepiye carane?
Amarga numeral Church
ditetepake minangka iterator
(fungsi kang ngetrapake
argumen kapisan ping $n$
marang argumen kapindho),
nilai wiwitan kita gawe
$\fn{true}$. Ing saben
langkah iterasi, asil
sing dibalekake
$\fn{false}$ tanpa
nggatekake asil langkah
sadurunge.
Langkah iki ditrapake
marang nilai wiwitan
ping $n$, lan asile
$\fn{true}$ yen lan
mung yen ora ana langkah
kang ditrapake, yaiku
nalika $n = 0$.

Saka perwakilan nilai
bebeneran iki, kita
uga bisa netepake
fungsi-fungsi bebeneran.
Iki rong conto, kanggo
negasi lan konjungsi:
\begin{align*}
  \fn{Not} & \ident \lambd[x][x\, \fn{false}\, \fn{true}]\\
  \fn{And} & \ident \lambd[x][\lambd[y][xy \,\fn{false}]]
\end{align*}
Fungsi $\fn{Not}$
nampa siji argumen lan
ngasilake $\fn{true}$
yen argumene $\fn{false}$,
lan $\fn{false}$
yen argumene~$\fn{true}$.
Fungsi $\fn{And}$
nampa rong nilai
bebeneran minangka
argumen, lan kudu
ngasilake $\fn{true}$
yen lan mung yen
loro-lorone
argumen~$\fn{true}$.
Nilai bebeneran iku
selektor kaya katrangan
ing ndhuwur, dadi yen
$x$ nilai bebeneran
lan ditrapake marang
rong argumen, asile
argumen kapisan yen
$x$ iku $\fn{true}$,
dene argumen kapindho
yen ora.
Saiki $\fn{And}$
nampa argumen $x$
lan $y$, banjur
menehake $y$ lan
$\fn{false}$ marang
argumen kapisan~$x$.
Yen $x$ nilai
bebeneran, asile
dievaluasi dadi~$y$
yen $x$ iku
$\fn{true}$, lan
dadi $\fn{false}$
yen $x$ iku
$\fn{false}$;
iki pancen asil kang
dikarepake.

Gatekna yen ing kene
kita nganggep mung
nilai bebeneran kang
dadi argumen
$\fn{And}$.
Yen fungsi iki
diwenehi suku liyane,
asile (yaiku wujud
normal, yen ana)
bisa wae dudu
nilai bebeneran.

\begin{prob}
  Tetepna fungsi
  $\fn{Or}$ lan
  $\fn{Xor}$ kang
  makili fungsi
  bebeneran disjungsi
  inklusif lan
  eksklusif nganggo
  pangodean nilai
  bebeneran minangka
  suku-$\lambd$.
\end{prob}

\end{document}
